\section{引言}

本学期课程作业三要求使用Vibe Coding——通过自然语言与AI编程工具交互来构建软件——制作一个个人主页。主页需包含三个区域：自我介绍（昵称/专业/年级+简介+代表图）、数据可视化（选一份关于自己的真实小数据，用合适的形式展示）、交互彩蛋（至少一个可交互触发的小设计）。评分从完整性、个性化、视觉与创意、Prompt工程和反思深度五个维度评估。

作为一门面向非计算机专业学生的课程，Vibe Coding的核心在于"人主导方向、AI执行实现"。这听起来理想，但实际操作中很快遇到一个根本问题：我的Prompt究竟该怎么说？模糊的描述如"帮我做个好看的页面"产出的结果毫无个性；而精确的设计描述又需要配色、字体、布局等方面的专业知识。对于一个既不会写代码、也没有设计背景的学生来说，如何在Vibe Coding中高效地驾驭AI，产出既有完整功能又有视觉品质的个人主页，是本实验试图回答的问题。

为此，我设计了一次自我对照实验：先使用Codex尝试"逐步构建"——让AI搭一个骨架，再一个模块一个模块往上加；然后改用Claude Code尝试"设计优先"——在写任何代码之前，先花大量时间与AI共同讨论和确定完整的设计方案。两种策略针对的是完全相同的任务目标。

本文基于这次实验，记录了两种策略各自的执行过程和结果，分析了增量构建失败的原因，总结了设计优先策略的关键环节，并给出了两种策略在实际效果上的对比。
